\documentclass[11pt]{amsart}

\usepackage{amsmath,amssymb,amsthm,mathtools}
\usepackage{booktabs}
\usepackage{hyperref}

\title{Solutions to Exercises of Fiber Bundlefrom Liang's Textbook: Differential Geometry and General Relativity}
\author{Haixuan Lin}
\date{\today}

\newtheorem{theorem}{Theorem}[section]
\newtheorem{proposition}[theorem]{Proposition}
\newtheorem{lemma}[theorem]{Lemma}
\newtheorem{corollary}[theorem]{Corollary}

\theoremstyle{definition}
\newtheorem{definition}[theorem]{Definition}
\newtheorem{example}[theorem]{Example}
\newtheorem{problem}[theorem]{Problem}

\theoremstyle{remark}
\newtheorem{remark}[theorem]{Remark}

\begin{document}

\begin{abstract}
  There is a well-known textbook in Chinese theoretic physics community called \emph{Differential Geometry and General Relativity} by Liang and Zhou \cite{liang2009differential} and in Volumn III of the series, there is a chapter on Fiber Bundle. This chapter is a good introduction to the concept of Fiber Bundle for us as beginners and is a good supplement to the book. Here I provide the solutions to the exercises in the chapter.
\end{abstract}

\maketitle

\section{Definitions and Theorems}

\begin{definition}
    \emph{Principal fiber bundle} consists of a \emph{bundle manifold} $P$, a \emph{base manifold} $M$ and a \emph{structure Lie group} $G$. They satisfy the following conditions:
    \begin{enumerate}
        \item $G$ acts freely on $P$ from the right via the action $R: P \times G \to P$;
        \item There exists a smooth projection map $\pi: P \to M$ such that for $\forall p \in P$, the fiber over $\pi(p)$ is given by $\pi^{-1}\left[\pi(p)\right] = \left\{ p g \mid g \in G \right\}$;
        \item $\forall x \in M$, there exists an open neighborhood $U$ of $x$ in $M$ and a diffeomorphism $T_U:\pi ^{-1}\left[ U \right] \rightarrow U\times G$, $T_U\left( p \right) =\left( \pi \left( p \right) ,S_U\left( p \right) \right)$, $\forall p\in \pi ^{-1}\left[ U \right]$, where $S_U:\pi ^{-1}\left[ U \right] \rightarrow G$ satisfies $S_U\left( pg \right) =S_U\left( p \right) g$, $\forall g\in G$.
    \end{enumerate}
\end{definition}


\begin{definition}
    Let $R_p:G\rightarrow P$, $R_p\left( g \right) \coloneq R_g\left( p \right)$, then we can define a fundamental vector field 
\end{definition}

The definition of a \emph{connection} on a principal fiber bundle has three equivalent forms:
\begin{definition}
    A connection on a principal fiber bundle $P\left(M,G\right)$ assigns to each point $p \in P$ a \emph{horizontal subspace} $H_p\subseteq T_pP$ such that:
    \begin{enumerate}
        \item $T_pP=V_p\oplus H_p$, $\forall p\in P$, where $V_p\coloneqq \left\{ X\in T_pP\mid \pi _*\left( X \right) =0 \right\}$ is \emph{vertical subspace;}
        \item $R_{g*}\left[ H_p \right] =H_{pg}$, $\forall p\in P$, $\forall g\in G$;
        \item $H_p$ is $C^\infty$ with respect to $p$.
    \end{enumerate}
\end{definition}
\begin{definition}
    A connection on a principal fiber bundle $P\left(M,G\right)$ is a $C^\infty$ $\mathfrak{g}$-valued 1-form on $P$, denoted by $\tilde{{\omega}}$, such that:
    \begin{enumerate}
        \item $\tilde{\omega}_p\left( A_{p}^{*} \right) =A$, $\forall A\in \mathfrak{g}$, $\forall p\in P$;
    \end{enumerate}
\end{definition}





\section{Problem 1}

\begin{problem}
  Prove Eq.(I-1-4), namely $S_U:\pi ^{=1}\left[ x \right] \rightarrow G$ and $R_{\breve{p}_U}:G\rightarrow \pi ^{-1}\left[ x \right]$ are a pair of inverse mappings.
\end{problem}
\begin{proof}
  
\end{proof}

\section{Problem 2}

\begin{problem}
  Prove Eq.(I-1-5), namely $R_g\circ R_p=R_{pf}$ for $\forall p\in P$, $\forall g\in G$.
\end{problem}
\begin{proof}
  
\end{proof}


\nocite{*}
\bibliographystyle{amsplain}
\bibliography{references}

\end{document}